\uuid{F1nA}
\titre{Logique}
\auteur{bodin ; croizat ; sacre}
\organisation{exo7}
\licence{CC BY-NC-SA 4.0}
\datecreate{2022-03-08}
\niveau{L1}
\module{Algèbre}
\chapitre{Logique, ensemble, raisonnement}
\sousChapitre{Logique}
\difficulte{2}
\theme{Logique}
\exo7qcmid{hilisit.logique.010}

\contenu{
\begin{question}
\qclassification{100.01}
Quelles sont les assertions vraies, quel que soit $x>0$, un réel strictement positif ?
\begin{answers}
	\bad{$\exists y > 0 \quad \ln(x) = y$}
	\good{$\exists y >0 \quad e^x = y$}
	\good{$\exists y > 0 \quad \ln(y) = x$}
	\bad{$\exists y > 0 \quad e^y = x$}
\end{answers}
\begin{explanations}
La valeur $y=\ln(x)$ est un nombre réel bien défini (car $x>0$), mais elle n'est pas nécessairement strictement supérieure à $0$ (prenez $x\leq 1$ pour que ce soit faux).\\
Tout nombre réel $x>0$ possède un antécédent par la fonction exponentielle. Si l'on note $y$ cet antécédent, cela signifie que $e^y = x$. Mais cet antécédent n'est pas nécessairement strictement positif ! En effet, pour $x=1$ par exemple, on aura $e^0 = 1$.\\
En revanche, le nombre $y=e^x$ sera toujours un réel strictement positif. Et en appliquant le logarithme à cette égalité, on obtient bien $\ln(y) = x$ avec ce même nombre $y>0$.
\end{explanations}
\end{question}
}
